% Generated draft: requires mathematical and source review.
\documentclass[11pt]{article}
\usepackage[margin=25mm]{geometry}
\usepackage{fontspec}
\setmainfont{Latin Modern Roman}
\usepackage{amsmath,amssymb,mathtools,unicode-math}
\setmathfont{Latin Modern Math}
\usepackage{enumitem,needspace,xurl,hyperref}
\setlength{\parindent}{0pt}
\setlength{\parskip}{5pt}
\setlength{\emergencystretch}{3em}
\newcommand{\N}{\mathbb{N}}
\newcommand{\Z}{\mathbb{Z}}
\newcommand{\Q}{\mathbb{Q}}
\newcommand{\R}{\mathbb{R}}
\newcommand{\C}{\mathbb{C}}
\newcommand{\F}{\mathbb{F}}
\newcommand{\A}{\mathbb{A}}
\newcommand{\PP}{\mathbb{P}}
\newcommand{\E}{\mathbb{E}}
\newcommand{\eps}{\varepsilon}
\newcommand{\dd}{\,\mathrm{d}}
\begin{document}
\section*{UnsolvedMath ID: 800592}
\textbf{Source problem code:} catalogue entry
\subsection{Definitions and mathematical statement}
Work with sets of reals equipped with their subspace topology. For such a set \(X\), let \(\mathcal O(X)\) denote its open covers, with the convention that no member of a cover equals \(X\). Let \(\Gamma(X)\) denote its infinite open covers \(\mathcal U\) such that each point of \(X\) belongs to all but finitely many members of \(\mathcal U\).

For classes of covers \(\mathcal A(X)\) and \(\mathcal B(X)\), the property \(\mathsf S_{\mathrm{fin}}(\mathcal A,\mathcal B)\) means that for every sequence \((\mathcal U_n)_{n\in\N}\) from \(\mathcal A(X)\), there are finite, possibly empty, families \(\mathcal V_n\subseteq\mathcal U_n\) such that \(\bigcup_{n\in\N}\mathcal V_n\in\mathcal B(X)\). The property \(\mathsf S_1(\mathcal A,\mathcal B)\) means that for every such sequence there are choices \(U_n\in\mathcal U_n\) such that \(\{U_n:n\in\N\}\in\mathcal B(X)\). In particular, \(\mathsf S_{\mathrm{fin}}(\mathcal O,\mathcal O)\) is the Menger property. Selection properties in an additivity expression denote the corresponding families of subsets of \(\R\).

For a family \(\mathcal J\subseteq\mathcal P(\R)\) for which a witnessing family exists, define
\[
\operatorname{add}(\mathcal J)=\min\left\{|\mathcal F|:\mathcal F\subseteq\mathcal J\text{ and }\bigcup\mathcal F\notin\mathcal J\right\}.
\]
The supplied definition instead writes \(\mathcal F\subseteq\mathcal I\), although \(\mathcal I\) is not defined. The displayed definition explicitly adopts the apparent intended correction \(\mathcal I=\mathcal J\); this requires source review.

Write \(f\leq^*g\) for \(f,g\in\N^{\N}\) when \(f(n)\leq g(n)\) for all but finitely many \(n\). The bounding number \(\mathfrak b\) is the least size of a family in \(\N^{\N}\) with no common \(\leq^*\)-upper bound. The dominating number \(\mathfrak d\) is the least size of a family \(D\subseteq\N^{\N}\) such that every \(f\in\N^{\N}\) satisfies \(f\leq^*g\) for some \(g\in D\). The notation \(\operatorname{cf}(\mathfrak d)\) denotes the cofinality of this cardinal.

For infinite subsets of \(\N\), write \(A\subseteq^*B\) when \(A\setminus B\) is finite. A family \(G\subseteq[\N]^{\aleph_0}\) is groupwise dense if it is closed under taking infinite almost subsets and, for every partition \((I_n)_{n\in\N}\) of \(\N\) into consecutive nonempty finite intervals, there is an infinite \(a\subseteq\N\) such that \(\bigcup_{n\in a}I_n\in G\). The groupwise density number \(\mathfrak g\) is the least cardinality of a collection of groupwise dense families with empty intersection.

For the contextual bounds below, \(\mathfrak h\) denotes the least size of a family of dense open subsets of the partially ordered set \(([\N]^{\aleph_0},\subseteq^*)\) whose intersection is not dense. Here open means downward closed in this order, and dense means that every condition has a stronger condition in the set. Let \(\mathcal N\) denote the ideal of Lebesgue null subsets of \(\R\).

The supplied source states the following contextual bounds, attributing them to the reference label \([\mathrm{AddQuad}]\):
\[
\max\{\mathfrak b,\mathfrak g\}\leq
\operatorname{add}(\mathsf S_{\mathrm{fin}}(\mathcal O,\mathcal O))
\leq\operatorname{cf}(\mathfrak d),
\]
\[
\mathfrak h\leq\operatorname{add}(\mathsf S_1(\Gamma,\Gamma))\leq\mathfrak b,
\qquad
\operatorname{add}(\mathcal N)\leq\operatorname{add}(\mathsf S_1(\mathcal O,\mathcal O)).
\]
These bounds are supplied background claims, not independently verified assertions in this draft.

The question attributed to Tsaban's 2006 list, Problem 18, is whether
\[
\operatorname{add}(\mathsf S_{\mathrm{fin}}(\mathcal O,\mathcal O))
=\max\{\mathfrak b,\mathfrak g\}.
\]
The record supplies no additional set-theoretic hypothesis. The cover conventions and cardinal-characteristic definitions above make explicit standard interpretations of notation that the extracted record does not itself define; their agreement with the full source requires review.
\subsection{Short English statement}
Does the least number of Menger subsets of \(\R\) whose union fails to be Menger equal \(\max\{\mathfrak b,\mathfrak g\}\)? The supplied registry labels this question open. The draft explicitly corrects an undefined symbol in the supplied additivity definition and records standard interpretations requiring source review.
\subsection{Sources}
\begin{itemize}
\item Local UnsolvedMath registry, permanent ID \(800592\), source problem code \texttt{catalogue entry}, supplied complete source record. The record lists Tsaban as proposer, gives the year 2006, and labels the problem open; this status is supplied metadata, not an independently verified current assessment.
\item Tsaban, \emph{Selection Principles and Special Sets of Reals: Open Problems} (2006), Problem 18. Supplied source URL: \url{https://arxiv.org/abs/math/0609601}. The record reports access on 2026-10-16 and extraction method \texttt{arxiv\_latex}; this draft does not claim independent access or source verification.
\item The supplied background, title, and contextual inequalities cite the label \([\mathrm{AddQuad}]\). No bibliographic expansion or separate URL for that reference is supplied, so its authorship and publication details remain unresolved.
\end{itemize}
\end{document}
